\appendices
\section{Appendix: Python Implementation Details}
The analysis was implemented as a Python desktop application with a GUI entry point in \texttt{main.py} and the main workflow class in \texttt{heartvolume/gui/app.py}.

\subsection{Program Structure}
The implementation is organized into three layers:
\begin{itemize}
    \item user interface and workflow orchestration (\texttt{HeartVolumeApp}),
    \item image/video interaction tools (frame selection, scale detection, distance measurement),
    \item cardiac metric computations in \texttt{heartvolume/core/calculations.py}.
\end{itemize}

\subsection{Core Computation Functions}
\begin{verbatim}
def calculate_volume_simple(L_cm, D_cm):
    return (pi / 6.0) * D_cm**2 * L_cm

def calculate_volume_simpson(L_cm, D_apex_cm, D_pm_cm, D_mv_cm):
    A_apex = pi * (D_apex_cm / 2.0)**2
    A_pm   = pi * (D_pm_cm   / 2.0)**2
    A_mv   = pi * (D_mv_cm   / 2.0)**2
    return (L_cm / 3.0) * (A_apex + A_mv + 4.0 * A_pm) / 3.0

def build_volume_result(EDV, ESV, HR=None):
    SV = EDV - ESV
    EF = 100.0 * SV / EDV
    CO = (SV * HR / 1000.0) if HR is not None else None
    return EDV, ESV, SV, EF, CO
\end{verbatim}

\subsection{Execution Workflow}
The software exposes three run modes corresponding to the report methods:
\begin{itemize}
    \item \texttt{run\_simple} for geometric EDV/ESV estimation,
    \item \texttt{run\_simpson} for 4-view modified Simpson analysis,
    \item \texttt{run\_doppler} for outflow-based SV and CO estimation.
\end{itemize}
Each run logs frame indices, calibrated scale values, measured distances, and final cardiac metrics used in the Results section.

\subsection{Reproducibility Notes}
To reproduce the computations, the final submission archives the exact input videos/images, selected ED/ES frames, calibrated cm/pixel values, and the console log generated by the application for each method.
